% ============================================================================
% BOLUM 2: PROBLEM TANIMI
% ============================================================================
% Kurumsal Temizlik Takip ve Yonetim Sistemi (KTTYS)
% Muhammet Anil YAGIZ - 213255046
% Kirikkale Universitesi - Bilgisayar Muhendisligi Bolumu
% ============================================================================

\chapter{PROBLEM TANIMI}
\label{chap:problem-tanimi}

% ============================================================================
% 2.1 GIRIS
% ============================================================================
\section{Giris}
\label{sec:problem-giris}

Kurumsal temizlik hizmetlerinin etkin yonetimi, ozellikle buyuk olcekli 
kurumlarda operasyonel verimlilik icin kritik oneme sahiptir. Universiteler, 
hastaneler, alisveris merkezleri ve kamu binalari gibi yogun insan 
trafigine sahip yapilarda temizlik hizmetlerinin koordinasyonu, takibi 
ve kalite kontrolu kritik bir yonetim fonksiyonu haline gelmektedir.

Bu bolumde, Kirikkale Universitesi ozelinde temizlik yonetimi 
sureclerinde karsilasilan problemler detayli olarak analiz 
edilmektedir.

% ============================================================================
% 2.2 MEVCUT DURUM ANALIZI
% ============================================================================
\section{Mevcut Durum Analizi}
\label{sec:mevcut-durum}

Kirikkale Universitesi, Turkiye'nin Ic Anadolu Bolgesinde yer alan 
ve cok sayida fakulte, yuksekokul, enstitu ve idari birimden olusan 
kapsamli bir yuksekogretim kurumudur.

\subsection{Fiziksel Altyapi Ozellikleri}
\label{subsec:fiziksel-altyapi}

Kirikkale Universitesi kampusunun fiziksel ozellikleri Tablo 
\ref{tab:fiziksel-altyapi}'da ozetlenmektedir.

\begin{table}[H]
\centering
\caption{Kirikkale Universitesi Fiziksel Altyapi Ozeti}
\label{tab:fiziksel-altyapi}
\begin{tabular}{|l|r|}
\hline
\textbf{Ozellik} & \textbf{Deger} \\
\hline
Toplam Fakulte Sayisi & 12 \\
\hline
Toplam Bina Sayisi & 45+ \\
\hline
Toplam Kapali Alan (m2) & 250.000+ \\
\hline
Derslik Sayisi & 300+ \\
\hline
Laboratuvar Sayisi & 150+ \\
\hline
Idari Ofis Sayisi & 200+ \\
\hline
Ortak Alan (Koridor, Hol) & 100+ \\
\hline
Tuvalet/Lavabo Alani & 200+ \\
\hline
Ogrenci Sayisi & 25.000+ \\
\hline
Akademik Personel & 1.500+ \\
\hline
Temizlik Personeli & 150+ \\
\hline
\end{tabular}
\end{table}

\subsection{Mevcut Temizlik Yonetim Sureci}
\label{subsec:mevcut-surec}

Kirikkale Universitesinde mevcut temizlik yonetim sureci, 
geleneksel yontemlere dayali olarak yurutulmektedir. Bu surec, 
Sekil \ref{fig:mevcut-surec}'te akis diyagrami olarak 
gosterilmektedir.

\begin{figure}[H]
\centering
\begin{tikzpicture}[
    node distance=1.2cm,
    startstop/.style={rectangle, rounded corners, minimum width=3cm, 
        minimum height=0.8cm, text centered, draw=black, fill=red!30},
    process/.style={rectangle, minimum width=3cm, minimum height=0.8cm, 
        text centered, draw=black, fill=orange!30},
    decision/.style={diamond, minimum width=2cm, minimum height=0.8cm, 
        text centered, draw=black, fill=green!30, aspect=2},
    io/.style={trapezium, trapezium left angle=70, trapezium right angle=110,
        minimum width=2cm, minimum height=0.8cm, text centered, 
        draw=black, fill=blue!30},
    arrow/.style={thick,->,>=stealth}
]
    \node (start) [startstop] {Baslangic};
    \node (plan) [process, below=of start] {Manuel Program Hazirlama};
    \node (assign) [process, below=of plan] {Sozel Gorev Atama};
    \node (execute) [process, below=of assign] {Temizlik Islemi};
    \node (check) [decision, below=of execute] {Kontrol?};
    \node (approve) [process, below left=1cm and 1cm of check] {Onay};
    \node (reject) [process, below right=1cm and 1cm of check] {Ret/Tekrar};
    \node (record) [io, below=of approve] {Kagit Kayit};
    \node (end) [startstop, below=of record] {Bitis};
    
    \draw [arrow] (start) -- (plan);
    \draw [arrow] (plan) -- (assign);
    \draw [arrow] (assign) -- (execute);
    \draw [arrow] (execute) -- (check);
    \draw [arrow] (check) -| node[pos=0.25, above] {Evet} (approve);
    \draw [arrow] (check) -| node[pos=0.25, above] {Hayir} (reject);
    \draw [arrow] (approve) -- (record);
    \draw [arrow] (reject) |- (execute);
    \draw [arrow] (record) -- (end);
\end{tikzpicture}
\caption{Mevcut Temizlik Yonetim Sureci Akis Diyagrami}
\label{fig:mevcut-surec}
\end{figure}

Mevcut surecin temel asamalari su sekilde ozetlenebilir:

\begin{enumerate}
    \item \textbf{Program Hazirlama:} Fakulte sekreterlikleri veya 
    amirler tarafindan haftalik/aylik temizlik programlari Excel 
    veya kagit uzerinde hazirlanmaktadir.
    
    \item \textbf{Gorev Atama:} Gorevler, sabah toplantilarinda veya 
    telefonla personele sozel olarak iletilmektedir.
    
    \item \textbf{Uygulama:} Temizlik personeli atanan gorevleri 
    kendi basina yurutmektedir.
    
    \item \textbf{Kontrol:} Denetleyiciler rastgele kontroller 
    yapmakta, ancak sistematik bir takip bulunmamaktadir.
    
    \item \textbf{Kayit:} Yapilan isler kagit formlara elle 
    yazilmakta veya hic kayit tutulmamaktadir.
\end{enumerate}

% ============================================================================
% 2.3 TESPIT EDILEN PROBLEMLER
% ============================================================================
\section{Tespit Edilen Problemler}
\label{sec:problemler}

Mevcut sistemin detayli analizi sonucunda, asagidaki temel 
problem kategorileri tespit edilmistir.

\subsection{Izlenebilirlik Problemleri}
\label{subsec:izlenebilirlik-problemleri}

\begin{enumerate}
    \item \textbf{Gorev Takibi Eksikligi:} Hangi lokasyonun ne zaman, 
    kim tarafindan temizlendiginin takibi mumkun degildir.
    
    \item \textbf{Kayit Guvenirliligi:} Kagit tabanli kayitlar 
    kolaylikla manipule edilebilmekte, kaybolabilmekte veya 
    zarar gorebilmektedir.
    
    \item \textbf{Gecmis Veri Erisimi:} Gecmis donemlere ait temizlik 
    verilerine erisim oldukca guctur ve zaman alicidir.
    
    \item \textbf{Anlik Durum Bilgisi:} Herhangi bir lokasyonun mevcut 
    temizlik durumu hakkinda anlik bilgi edinmek mumkun degildir.
\end{enumerate}

\subsection{Koordinasyon Problemleri}
\label{subsec:koordinasyon-problemleri}

\begin{enumerate}
    \item \textbf{Iletisim Kopukluklari:} Yonetim, denetleyiciler ve 
    temizlik personeli arasindaki iletisim cogunlukla anlik ve 
    yapilandirilmamis bir sekilde gerceklesmektedir.
    
    \item \textbf{Program Cakismalari:} Farkli birimlerdeki programlar 
    arasinda koordinasyon eksikligi nedeniyle cakismalar 
    yasanabilmektedir.
    
    \item \textbf{Acil Durum Mudahalesi:} Beklenmeyen temizlik 
    ihtiyaclarina (dokuntu, kaza vb.) hizli mudahale koordinasyonu 
    zorlasmaktadir.
    
    \item \textbf{Personel Ikamesi:} Izinli veya hasta personelin 
    yerine ikame atama sureci verimsiz calismaktadir.
\end{enumerate}

\subsection{Kaynak Yonetimi Problemleri}
\label{subsec:kaynak-problemleri}

\begin{enumerate}
    \item \textbf{Is Yuku Dengesizligi:} Bazi personele asiri yuklenme 
    olurken, bazilari dusuk kapasiteyle calismaktadir.
    
    \item \textbf{Zaman Optimizasyonu:} Personelin bina icindeki 
    hareket rotalari optimize edilmemekte, gereksiz zaman kaybi 
    yasanmaktadir.
    
    \item \textbf{Malzeme Takibi:} Temizlik malzemelerinin kullanimi 
    ve stok durumu etkin bir sekilde takip edilememektedir.
    
    \item \textbf{Ekipman Bakimi:} Temizlik ekipmanlarinin bakim 
    zamanlarinin takibi yapilamamaktadir.
\end{enumerate}

\subsection{Raporlama ve Analiz Problemleri}
\label{subsec:raporlama-problemleri}

\begin{enumerate}
    \item \textbf{Performans Olcumu:} Bireysel ve takim performansinin 
    objektif olcumu mumkun degildir.
    
    \item \textbf{Istatistiksel Analiz:} Temizlik hizmetlerine iliskin 
    trend analizleri, karsilastirmali degerlendirmeler 
    yapilamamaktadir.
    
    \item \textbf{Karar Destek:} Yoneticiler, veri destekli kararlar 
    almak icin gerekli bilgiye erisememektedir.
    
    \item \textbf{Maliyet Analizi:} Birim basina temizlik maliyetinin 
    hesaplanmasi gucleşmektedir.
\end{enumerate}

% ============================================================================
% 2.4 PROBLEM AGACI ANALIZI
% ============================================================================
\section{Problem Agaci Analizi}
\label{sec:problem-agaci}

Tespit edilen problemlerin sistematik analizi icin problem agaci 
(problem tree) yontemi uygulanmistir. Bu yontem, temel problemin 
nedenlerini ve sonuclarini hiyerarsik olarak gorsellestirmektedir.

\begin{figure}[H]
\centering
\begin{tikzpicture}[
    node distance=1.2cm,
    box/.style={rectangle, draw, text width=3.2cm, text centered, 
                minimum height=0.8cm, font=\small},
    mainbox/.style={rectangle, draw, fill=red!30, text width=5cm, 
                    text centered, minimum height=1cm, font=\bfseries},
    causebox/.style={rectangle, draw, fill=orange!30, text width=2.8cm, 
                     text centered, minimum height=0.7cm, font=\small},
    effectbox/.style={rectangle, draw, fill=blue!30, text width=2.8cm, 
                      text centered, minimum height=0.7cm, font=\small}
]
    % Ana Problem
    \node[mainbox] (main) {Temizlik Hizmetlerinin Verimsiz Yonetimi};
    
    % Nedenler (alt)
    \node[causebox, below left=1.5cm and 3cm of main] (cause1) 
        {Dijital Altyapi Eksikligi};
    \node[causebox, below left=1.5cm and 0.5cm of main] (cause2) 
        {Manuel Surecler};
    \node[causebox, below right=1.5cm and 0.5cm of main] (cause3) 
        {Iletisim Yetersizligi};
    \node[causebox, below right=1.5cm and 3cm of main] (cause4) 
        {Egitim Eksikligi};
    
    % Sonuclar (ust)
    \node[effectbox, above left=1.5cm and 3cm of main] (effect1) 
        {Dusuk Verimlilik};
    \node[effectbox, above left=1.5cm and 0.5cm of main] (effect2) 
        {Maliyet Artisi};
    \node[effectbox, above right=1.5cm and 0.5cm of main] (effect3) 
        {Kalite Dususu};
    \node[effectbox, above right=1.5cm and 3cm of main] (effect4) 
        {Paydas Memnuniyetsizligi};
    
    % Oklar
    \draw[->] (cause1) -- (main);
    \draw[->] (cause2) -- (main);
    \draw[->] (cause3) -- (main);
    \draw[->] (cause4) -- (main);
    \draw[->] (main) -- (effect1);
    \draw[->] (main) -- (effect2);
    \draw[->] (main) -- (effect3);
    \draw[->] (main) -- (effect4);
    
    % Etiketler
    \node[left=0.3cm of cause1, font=\small\itshape] {Nedenler};
    \node[left=0.3cm of effect1, font=\small\itshape] {Sonuclar};
    
\end{tikzpicture}
\caption{Temizlik Yonetimi Problem Agaci}
\label{fig:problem-agaci}
\end{figure}

\subsection{Kok Nedenler}
\label{subsec:kok-nedenler}

Problem agaci analizinde tespit edilen kok nedenler sunlardir:

\begin{description}
    \item[Dijital Altyapi Eksikligi:] Universitede temizlik 
    hizmetlerinin yonetimi icin ozel olarak tasarlanmis bir 
    yazilim sistemi bulunmamaktadir.
    
    \item[Manuel Surecler:] Planlama, atama, takip ve raporlama 
    sureclerinin tamami manuel olarak yurutulmektedir.
    
    \item[Iletisim Yetersizligi:] Hiyerarsik yapi icindeki farkli 
    kademeler arasinda etkin bir iletisim kanali bulunmamaktadir.
    
    \item[Egitim Eksikligi:] Personelin dijital araclarin kullanimi 
    konusunda yeterli egitimi bulunmamaktadir.
\end{description}

\subsection{Sonuclar ve Etkileri}
\label{subsec:sonuclar}

Kok nedenlerin yarattigi olumsuz sonuclar sunlardir:

\begin{description}
    \item[Dusuk Verimlilik:] Personel zamaninin onemli bir kismi 
    koordinasyon ve iletis sorunlarina harcanmaktadir.
    
    \item[Maliyet Artisi:] Verimsiz kaynak kullanimi, fazla mesai 
    ihtiyaci ve malzeme israfi toplam maliyetleri artirmaktadir.
    
    \item[Kalite Dususu:] Denetim eksikligi ve takip yetersizligi, 
    temizlik hizmetlerinin kalitesini olumsuz etkilemektedir.
    
    \item[Paydas Memnuniyetsizligi:] Ogrenciler, akademik personel 
    ve ziyaretciler temizlik hizmetlerinden memnun kalmayabilmektedir.
\end{description}

% ============================================================================
% 2.5 PAYDAS ANALIZI
% ============================================================================
\section{Paydas Analizi}
\label{sec:paydas-analizi}

Sistemin basarili bir sekilde tasarlanmasi ve uygulanmasi icin 
tum paydaslarin tanimlanmasi ve ihtiyaclarinin analiz edilmesi 
gerekmektedir.

\subsection{Birincil Paydaslar}
\label{subsec:birincil-paydaslar}

\begin{table}[H]
\centering
\caption{Birincil Paydas Analizi}
\label{tab:birincil-paydaslar}
\begin{tabular}{|p{2.5cm}|p{4cm}|p{5cm}|}
\hline
\textbf{Paydas} & \textbf{Rol} & \textbf{Ihtiyaclar/Beklentiler} \\
\hline
Temizlik Personeli & Temizlik gorevlerini yurutur & Acik gorev tanimlari, 
adil is dagilimi, kolay kullanim \\
\hline
Denetleyiciler & Personeli denetler ve gorevleri onaylar & 
Gercek zamanli takip, onay mekanizmasi, raporlama \\
\hline
Fakulte Yoneticileri & Fakulte bazinda temizlik yonetimi & Program 
olusturma, personel atama, analitik gorunum \\
\hline
Genel Yonetim & Sistem geneli yonetim & Merkezi kontrol, 
kapsamli raporlar, sistem konfigurasyonu \\
\hline
\end{tabular}
\end{table}

\subsection{Ikincil Paydaslar}
\label{subsec:ikincil-paydaslar}

\begin{itemize}
    \item \textbf{Ogrenciler:} Temiz ve hijyenik ortamlarda 
    egitim gormek istemektedirler.
    
    \item \textbf{Akademik Personel:} Temiz calisma ortami 
    beklentisi icindedirler.
    
    \item \textbf{Ziyaretciler:} Kampusun temizlik durumu 
    kurumsal imaji etkilemektedir.
    
    \item \textbf{Idari Personel:} Ortak alanlarin temizligi 
    calisma verimliligini etkilemektedir.
    
    \item \textbf{Malzeme Tedarikcileri:} Temizlik malzemesi 
    tuketim verilerine ihtiyac duyulmaktadir.
\end{itemize}

% ============================================================================
% 2.6 COZUM HEDEFLERI
% ============================================================================
\section{Cozum Hedefleri}
\label{sec:cozum-hedefleri}

Tespit edilen problemlerin cozumune yonelik olarak gelistirilen 
sistemin karsilamasi gereken hedefler asagida tanimlanmaktadir.

\subsection{Birincil Hedefler}
\label{subsec:birincil-hedefler}

\begin{enumerate}
    \item \textbf{Dijital Donusum:} Manuel temizlik yonetim 
    sureclerinin tamamen dijitallestirilmesi.
    
    \item \textbf{Merkezi Yonetim:} Tum fakultelerin tek bir 
    platform uzerinden yonetilebilmesi.
    
    \item \textbf{Izlenebilirlik:} Tum gorevlerin ve islemlerin 
    kayit altina alinarak izlenebilir hale getirilmesi.
    
    \item \textbf{Verimlilik Artisi:} Koordinasyon ve iletisim 
    sureclerinin optimizasyonu ile verimlilik artisi saglanmasi.
\end{enumerate}

\subsection{Ikincil Hedefler}
\label{subsec:ikincil-hedefler}

\begin{enumerate}
    \item \textbf{Kullanici Deneyimi:} Tum kullanici gruplari icin 
    kolay ve sezgisel bir arayuz sunulmasi.
    
    \item \textbf{Mobil Erisim:} Sistemin mobil cihazlardan 
    erisilebilir olmasi.
    
    \item \textbf{Raporlama Kapasitesi:} Kapsamli ve ozellestirilebilir 
    raporlar olusturulabilmesi.
    
    \item \textbf{Olceklenebilirlik:} Sistemin gelecekte universitenin 
    tum birimlerini kapsayacak sekilde genisletilebilmesi.
\end{enumerate}

% ============================================================================
% 2.7 BASARI KRITERLERI
% ============================================================================
\section{Basari Kriterleri}
\label{sec:basari-kriterleri}

Projenin basarisinin degerlendirilmesi icin belirlenen olculebilir 
kriterler Tablo \ref{tab:basari-kriterleri}'de sunulmaktadir.

\begin{table}[H]
\centering
\caption{Proje Basari Kriterleri}
\label{tab:basari-kriterleri}
\begin{tabular}{|p{4cm}|p{3cm}|p{4cm}|}
\hline
\textbf{Kriter} & \textbf{Hedef Deger} & \textbf{Olcum Yontemi} \\
\hline
Gorev tamamlama orani & 95\%+ & Sistem kayitlari \\
\hline
Ortalama gorev onay suresi & $<$ 2 saat & Zaman damgasi analizi \\
\hline
Sistem kullanim orani & 90\%+ & Aktif kullanici takibi \\
\hline
Kullanici memnuniyeti & 4/5+ & Anket calismasi \\
\hline
Raporlama suresi & $<$ 1 dakika & Performans testi \\
\hline
Sistem uptime & 99\%+ & Sunucu izleme \\
\hline
\end{tabular}
\end{table}

Bu bolumde sunulan problem analizi, sistemin tasarim ve gelistirme 
sureclerine temel teskil etmektedir. Bir sonraki bolumde, bu alanda 
yapilmis benzer calismalar incelenecektir.
